\documentclass[12pt,a4paper]{article}

% Required Packages
\usepackage[utf8]{inputenc}
\usepackage{graphicx}
\usepackage{amsmath}
\usepackage{xcolor}
\usepackage{geometry}

% Page Margins
\geometry{margin=1in}

% Document Information
\title{\textbf{Introduction to LaTeX}}
\author{Abhishek Narayan Hathwala}
\date{\today}

\begin{document}

% Title Page
\maketitle

% Table of Contents
\tableofcontents
\newpage

% Section 1
\section{Introduction}

LaTeX is a powerful document preparation system used for
creating professional academic, scientific, and technical documents.

It provides excellent formatting capabilities and supports
mathematical equations, tables, figures, and references.

% Section 2
\section{Features of LaTeX}

Some important features of LaTeX are:

\begin{itemize}
    \item Professional document formatting.
    \item Mathematical equation support.
    \item Automatic numbering.
    \item Reference management.
    \item Structured document preparation.
\end{itemize}

% Section 3
\section{Document Structure}

A LaTeX document generally consists of two parts:

\begin{enumerate}
    \item Preamble
    \item Document Body
\end{enumerate}

The preamble contains document settings and packages,
while the document body contains the actual content.

% Section 4
\section{Text Formatting}

LaTeX supports different text formatting styles.

\begin{itemize}
    \item \textbf{Bold Text}
    \item \textit{Italic Text}
    \item \underline{Underlined Text}
    \item \texttt{Monospaced Text}
\end{itemize}

% Section 5
\section{Mathematical Expression}

LaTeX is widely used for mathematical expressions.

Example:

\[
a^2+b^2=c^2
\]

Another example:

\[
E=mc^2
\]

% Section 6
\section{Conclusion}

LaTeX is an essential tool for technical writing.
It helps create structured, consistent, and professional
documents with excellent formatting
\end{document}